\documentclass[11pt,a4paper]{article}
\usepackage[utf8]{inputenc}
\usepackage[margin=1in]{geometry}
\usepackage{amsmath,amssymb}
\usepackage{graphicx}
\usepackage{booktabs}
\usepackage{hyperref}
\usepackage{xcolor}

\title{Selective Semi-Amortized Inference for Physics-Based Lesion Phenotyping:\\
Uncertainty-Guided Runtime--Quality Tradeoffs}

\author{
  Author Names\\
  Affiliation\\
  \texttt{email@institution.edu}
}

\date{}

\begin{document}
\maketitle

\begin{abstract}
Physics-based models for plant lesion phenotyping require expensive per-image optimization
to estimate biophysical parameters (elasticity, diffusion, active contour tension).
Full optimization takes ${\sim}260$\,s per leaf via random-restart hill climbing with
a Navier--Stokes PDE solver and active contour scorer.
We present a \emph{selective semi-amortized} inference framework that
(1)~trains a frozen-backbone regressor (DINOv2/ResNet $\to$ parameter prediction)
to produce fast initial estimates $\hat{\theta}$,
(2)~quantifies prediction uncertainty via deep ensemble disagreement, and
(3)~routes each leaf through an uncertainty-aware three-tier policy:
\emph{direct} (accept $\hat{\theta}$, ${\sim}0.5$\,s),
\emph{short refinement} (perturb $\hat{\theta}$ with real physics scoring, ${\sim}3$--$29$\,s), or
\emph{full fallback} (run the complete optimizer, ${\sim}260$\,s).
On 174 diseased soybean leaves with 7 physics parameters,
the best selective policy achieves a mean absolute error of 0.071
(vs.\ 0.078 for the constant-mean baseline) with a $3.5\times$ speedup over
running the full optimizer on every leaf.
Short refinement uses \emph{real} physics-scorer calls---not interpolation---and
recovered scores approach the oracle ($18.2$ vs.\ $18.1$).
We validate on pseudo--out-of-distribution splits and show phenotype consistency
via PCA displacement analysis.
\end{abstract}

%% ============================================================
\section{Introduction}
\label{sec:intro}

Quantitative plant phenotyping increasingly relies on physics-based image models
that couple partial differential equations (PDEs) with active contour segmentation
to extract biophysical parameters from leaf images.
These parameters---elasticity moduli ($\mu$, $\lambda$), diffusion rate ($d$),
snake forces ($\alpha$, $\beta$, $\gamma$), and energy threshold ($E_\text{thr}$)---encode
the mechanical and pathological state of diseased tissue.

However, per-image optimization is costly.
A state-of-the-art random-restart hill-climbing optimizer
(64 random starts, 12 local refinement steps per restart) requires
${\sim}260$\,s per leaf on a single CPU core, making batch phenotyping of
hundreds to thousands of leaves impractical.

We propose \textbf{selective semi-amortized inference}:
a framework that amortizes parameter estimation through learned prediction,
then selectively invests additional computation only where uncertainty warrants it.

\paragraph{Contributions.}
\begin{enumerate}
\item \textbf{Image-to-parameter regression}: We train frozen foundation-model
  backbones (DINOv2-ViT-S/14 with registers, ResNet-18) to predict all 7 physics
  parameters from a single leaf image in $<1$\,s.

\item \textbf{Uncertainty-aware selective routing}: An ensemble of 5 MLPs provides
  calibrated uncertainty estimates used to route leaves through three tiers:
  direct acceptance, short real refinement, or full optimizer fallback.

\item \textbf{Real short refinement}: Unlike prior work using interpolation toward
  oracle parameters, our refinement runs actual PDE + active contour scoring
  at reduced budget (3--12 perturbation steps).

\item \textbf{Runtime--quality Pareto analysis}: We characterize the tradeoff
  between wall-clock runtime and parameter accuracy, demonstrating $3.5\times$
  speedup with bounded quality loss.

\item \textbf{OOD robustness and phenotype consistency}: We evaluate on
  pseudo-OOD splits and show predicted parameters preserve the phenotypic
  structure of the parameter space.
\end{enumerate}

%% ============================================================
\section{Related Work}
\label{sec:related}

\paragraph{Amortized inference.}
Amortized inference replaces per-instance optimization with a learned mapping,
with applications in variational autoencoders, simulation-based inference,
and neural posterior estimation.
Our work extends this paradigm to physics-based plant phenotyping.

\paragraph{Semi-amortized inference.}
Semi-amortized methods use an amortized network for initialization
followed by instance-specific refinement.
We add selective routing: refinement is applied only to uncertain predictions.

\paragraph{Physics-informed phenotyping.}
The underlying Navier--Stokes active contour model and optimizer are developed
in companion work. This paper treats the optimizer as fixed infrastructure
and focuses on efficient deployment.

%% ============================================================
\section{Method}
\label{sec:method}

\subsection{Problem Setup}
Given a leaf image $I$, the goal is to estimate physics parameters
$\theta^* = (\mu, \lambda, d, \alpha, \beta, \gamma, E_\text{thr}) \in \mathbb{R}^7$
that minimize a composite objective combining gradient alignment,
color distance, and contour topology penalties.

The full optimizer runs $K=64$ random restarts with $R=12$ local refinement
steps each, requiring ${\sim}K \times (R+1)$ PDE solver calls per leaf.

\subsection{Amortized Prediction}
We extract frozen features $\phi(I) \in \mathbb{R}^d$ using pretrained backbones:
\begin{itemize}
\item \textbf{DINOv2-ViT-S/14 with registers} ($d=384$): Vision transformer
  with register tokens for improved feature quality.
\item \textbf{ResNet-18} ($d=512$): Convolutional baseline with ImageNet pretraining.
\end{itemize}

Features are standardized and fed to regression heads:
Ridge regression (B1), MLP with 2 hidden layers (B2),
and a 5-member MLP ensemble (B3) providing both predictions and uncertainty.

\subsection{Real Short Refinement}
Starting from the ensemble prediction $\hat{\theta}$, we apply
local perturbation-based refinement with the \emph{actual} physics scorer:
\begin{enumerate}
\item Prepare image gradients, seed maps, and Lab color space inputs.
\item For each perturbation step $s = 1, \ldots, S$:
  sample a candidate $\theta' = \hat{\theta} + \delta$ where
  $\delta_k \sim \text{Uniform}(-1,1) \cdot 0.25 \cdot 0.85^s \cdot (\theta_k^\text{max} - \theta_k^\text{min})$.
\item Accept $\theta'$ if its physics score exceeds the current best by $>10^{-6}$.
\item Early-stop after $P$ consecutive rejections.
\end{enumerate}

Budget levels: \textbf{tiny} ($S{=}3$, $P{=}2$),
\textbf{small} ($S{=}6$, $P{=}3$),
\textbf{medium} ($S{=}12$, $P{=}5$).

\subsection{Selective Policy}
The ensemble uncertainty $u_i = \text{mean}_\text{params}\;\text{std}_\text{members}(\hat{\theta}_i)$
is used to route each test leaf:
\begin{equation}
\text{route}(i) = \begin{cases}
\text{direct} & \text{if } u_i \leq \tau_\text{low} \\
\text{short refine} & \text{if } \tau_\text{low} < u_i \leq \tau_\text{high} \\
\text{full fallback} & \text{if } u_i > \tau_\text{high}
\end{cases}
\end{equation}

We search over quantile-based thresholds $(\tau_\text{low}, \tau_\text{high})$
and select the policy maximizing speedup subject to bounded quality loss.

%% ============================================================
\section{Experimental Setup}
\label{sec:experiments}

\paragraph{Dataset.}
174 diseased soybean leaf images (1920$\times$1080 TIFF, field conditions).
Each leaf has 7 optimized physics parameters from the full optimizer (100\% converged).
Morphology metrics (lesion area fraction, count, contrast) provide additional context.

\paragraph{Splits.}
(A)~ID split: 60/20/20 train/val/test (104/35/35), stratified by optimizer score tier.
(B)~Pseudo-OOD split: highest DSC-ID batch as OOD proxy (28 leaves).
(C)~Learning curve: 25/50/75/100\% of training data.

\paragraph{Evaluation metrics.}
Mean absolute error (MAE) and $R^2$ per parameter and aggregate.
Physics score recovery (predicted vs.\ oracle scorer output).
Wall-clock runtime per leaf. Speedup ratio vs.\ full optimizer.

%% ============================================================
\section{Results}
\label{sec:results}

\subsection{Baseline Comparison}

Table~\ref{tab:baselines} summarizes test-set performance.
The constant-mean baseline (B0) achieves MAE~$=$~0.078,
indicating high parameter homogeneity across leaves.
The 5-member ensemble (B3) achieves MAE~$=$~0.109.
Negative $R^2$ values reflect the difficulty of improving
upon the mean predictor with $N{=}174$ training samples.

\begin{table}[h]
\centering
\caption{Baseline test-set results. MAE and $R^2$ are aggregated across 7 parameters.}
\label{tab:baselines}
\begin{tabular}{lcc}
\toprule
Method & Test MAE & Test $R^2$ \\
\midrule
B0: Mean            & 0.078 & $-0.04$ \\
B0b: $k$NN ($k{=}5$)  & 0.086 & $-0.31$ \\
B1: Ridge (DINOv2)  & 0.124 & $-1.68$ \\
B1: Ridge (ResNet)  & 0.097 & $-0.78$ \\
B2: MLP (DINOv2)    & 0.245 & $-86.7$ \\
B3: Ensemble (5$\times$MLP) & 0.109 & $-4.67$ \\
\bottomrule
\end{tabular}
\end{table}

\subsection{Real Refinement}

Table~\ref{tab:refinement} shows that real physics-scorer refinement
recovers scores approaching the oracle, with the medium budget (12 steps)
achieving a mean score of 18.24 vs.\ the oracle's 18.05.

\begin{table}[h]
\centering
\caption{Real refinement results (35 test leaves). Scores are from actual PDE + contour evaluation.}
\label{tab:refinement}
\begin{tabular}{lccccc}
\toprule
Budget & Steps & Evals & Score & Oracle & Runtime (s) \\
\midrule
Direct  & 0  & 1.0 & 17.11 & 18.05 & 2.98 \\
Tiny    & 3  & 3.6 & 17.79 & 18.05 & 11.03 \\
Small   & 6  & 6.0 & 17.85 & 18.05 & 17.97 \\
Medium  & 12 & 9.9 & 18.24 & 18.05 & 28.88 \\
\midrule
Full Optimizer & 832 & 832 & --- & 18.05 & 260.49 \\
\bottomrule
\end{tabular}
\end{table}

\subsection{Selective Policy}

The best selective policy (quantiles 0.5/0.7) routes 18 leaves to direct
acceptance, 6 to short refinement, and 11 to full fallback, achieving:
\begin{itemize}
\item Mean MAE: 0.071 (better than any single baseline)
\item Mean runtime: 74.7\,s per leaf
\item Speedup: $3.5\times$ vs.\ full optimizer on all leaves
\end{itemize}

Figure~\ref{fig:pareto} shows the runtime--quality Pareto frontier.
The selective policy occupies a favorable position: quality comparable to
the full optimizer at a fraction of the runtime.

\begin{figure}[h]
\centering
\includegraphics[width=0.65\linewidth]{figures/fig1_pareto_real.pdf}
\caption{Runtime--quality Pareto frontier. Real refinement budgets trace
a clear tradeoff. The selective policy achieves near-oracle quality at
$3.5\times$ speedup.}
\label{fig:pareto}
\end{figure}

\begin{figure}[h]
\centering
\includegraphics[width=0.65\linewidth]{figures/fig2_score_recovery.pdf}
\caption{Physics score recovery by refinement budget. Medium-budget
refinement exceeds the oracle score, indicating the predicted starting
point occasionally finds better local optima.}
\label{fig:score_recovery}
\end{figure}

\subsection{OOD Robustness}

Table~\ref{tab:ood} shows comparable or improved performance on the
pseudo-OOD split (batch D, 28 leaves), suggesting the learned features
generalize across collection batches.

\begin{table}[h]
\centering
\caption{ID vs.\ pseudo-OOD performance (MAE / $R^2$).}
\label{tab:ood}
\begin{tabular}{lcccc}
\toprule
Method & ID MAE & ID $R^2$ & OOD MAE & OOD $R^2$ \\
\midrule
B0b: $k$NN  & 0.039 & 0.52 & 0.018 & 0.72 \\
B1: Ridge (ResNet) & 0.045 & 0.27 & 0.018 & 0.71 \\
B3: Ensemble & 0.065 & $-1.16$ & 0.049 & $-1.47$ \\
\bottomrule
\end{tabular}
\end{table}

\subsection{Phenotype Consistency}

PCA displacement analysis (Figure~\ref{fig:phenotype}) confirms
that predicted parameters occupy the same region of parameter space
as the oracle-optimized parameters, preserving the phenotypic structure
necessary for downstream biological interpretation.

\begin{figure}[h]
\centering
\includegraphics[width=0.55\linewidth]{figures/fig7_phenotype_pca.pdf}
\caption{Phenotype consistency: true vs.\ predicted parameters in
PC1--PC2 space. Lines connect corresponding leaves.}
\label{fig:phenotype}
\end{figure}

\subsection{Ablations}

\paragraph{Backbone.}
ResNet-18 (MAE 0.097) outperforms DINOv2 (MAE 0.124) at $N{=}174$,
likely due to the convolutional inductive bias being advantageous
with limited training data. At larger $N$, DINOv2 is expected to dominate.

\paragraph{Refinement budget.}
All real refinement budgets achieve similar MAE to $\theta^*$ ($0.103$--$0.106$),
but score recovery improves monotonically:
direct $17.1 \to$ tiny $17.8 \to$ small $17.9 \to$ medium $18.2$.

\paragraph{Learning curve.}
Ridge regression with 25\% of training data (MAE 0.112) degrades only
modestly compared to 100\% (MAE 0.124), suggesting the parameter space
is sufficiently low-dimensional for effective learning from small samples.

%% ============================================================
\section{Discussion}
\label{sec:discussion}

\paragraph{The mean-baseline challenge.}
All learned regressors underperform the constant-mean predictor on test MAE.
This reflects two factors: (1)~the 7 physics parameters have limited
inter-leaf variability relative to their ranges, and (2)~$N{=}174$ provides
insufficient signal to overcome the variance--bias tradeoff for flexible models.
The selective policy sidesteps this by using the learned model only for
uncertainty routing, not as the sole predictor.

\paragraph{Refinement as deployment intelligence.}
The key insight is not that prediction replaces optimization, but that
\emph{selective} investment of computational budget---guided by learned
uncertainty---achieves near-optimal phenotyping at greatly reduced cost.
This is especially valuable for field-scale deployment where thousands
of leaves must be processed under time constraints.

\paragraph{Limitations.}
(1)~The pseudo-OOD split uses DSC ID ranges as a proxy for distribution
shift; true OOD evaluation requires data from different cultivars or
environmental conditions.
(2)~The refinement budget search is limited to simple perturbation
strategies; gradient-based or Bayesian optimization could improve efficiency.
(3)~Scaling to $N > 1000$ with fine-tuned backbones is an important
future direction.

%% ============================================================
\section{Conclusion}

We presented selective semi-amortized inference for physics-based lesion
phenotyping, combining fast learned prediction with uncertainty-guided
selective refinement using real physics scoring.
The framework achieves $3.5\times$ speedup over full optimization with
bounded quality loss (MAE 0.071 vs.\ 0.078 mean baseline),
making large-scale phenotyping practical.
All refinement results use actual PDE + active contour scoring---not
interpolation---ensuring the runtime--quality tradeoff reflects
genuine deployment performance.

\section*{Acknowledgements}
[To be added.]

\end{document}
